\documentclass[book.tex]{subfiles}
\begin{document}

\chapter{Hopfield Networks}
%\url{https://journals.aps.org/prl/abstract/10.1103/PhysRevLett.98.146401}
%\url{file:///Users/bharath/Downloads/behler2007.pdf}
\label{chap:behler_parinello}
\noindent\rule{11cm}{0.4pt} \\
\textbf{Prerequisites:} \autoref{chap:ml_basics} \\
\textbf{Difficulty Level:} ** \\
\noindent\rule{11cm}{0.4pt}
\vspace{5mm}

Hopfield networks are a type of neural network that is not trained with gradient descent. Instead, a neuroscience-inspired update, the Hebbian rule, is used to update weights. Informally, the rule used is "neurons that fire together wire together."
\end{document}